\subsection{An explicit finite cubic-centralizer detector}
\label{sec:finite-cubic-detector}

\begin{result}{A.4.3}{Proved}{A finite five-point quotient removes the cubic-frame ambiguity}
Put $P=\Gamma_{0,5}=P_4/\langle z_4\rangle$, with ordered generators
\[
 (a,b,x,y,z)=(A_{13},A_{23},A_{14},A_{24},A_{34}),
 \qquad z_4=(\sigma _1\sigma _2\sigma _3)^4.
\]
Let $\ell:P\to\mathbb F_2^5$ be abelianization modulo two in this
ordered basis.  Let $q_3:P\to\operatorname{SL}_2(\mathbb F_3)$ be
the restriction of the literal braid representation
\[
 \sigma _1,\sigma _3\longmapsto U=
 \begin{pmatrix}1&1\\0&1\end{pmatrix},\qquad
 \sigma _2\longmapsto L=
 \begin{pmatrix}1&0\\-1&1\end{pmatrix}.
\]
The map $q=(\ell,q_3)$ is onto
$Q=\mathbb F_2^5\times\operatorname{SL}_2(\mathbb F_3)$, of order
$768$.  Define the actual normal subgroup and finite quotient
\[
 N=\ker q,\qquad R=P/[N,N]N^3.
\]
Thus the kernel in this formula is the subgroup generated by the
commutators and cubes in $N$.  Its profinite version is their closed
normal subgroup; the two constructions give the same finite quotient.
One has
\[
 |R|=768\cdot3^{114}.
\]
The element $d=\sigma _1^2\sigma _2\sigma _3$ has order three in
$B_4/\langle z_4\rangle$ and normalizes $P$ and $N$.  For its induced
action on $R$,
\[
 \operatorname{im}\bigl(C_R(d)\longrightarrow Q\bigr)=1,
 \qquad |C_R(d)|=3^{40}.                         \tag{FCQ.1}
\]
In particular the residue
\[
 \ell(h)=(0,1,0,1,1),\qquad q_3(h)=I              \tag{FCQ.2}
\]
cannot be the image of an element of $C_R(d)$.  The vector in
\textup{(FCQ.2)} is the class of the cyclic boundary generator
$X_5=(\sigma _3\sigma _2\sigma _3)^2=b y z$ modulo two.  This is the
specific finite centralizer ambiguity isolated in DELTA A.112.
\end{result}

\begin{proof}
We give the exact finite construction, the complete elimination
criterion, and the certificate for every candidate residue.  All
matrix arithmetic below takes place in $\mathbb F_3$ unless another
ring is specified.

\paragraph{The presentation and the finite cover.}
The free-by-free presentation of $P$ has the following six relators,
where $c(u,v)=v^{-1}u^{-1}vu$:
\[
\begin{array}{lll}
 axa^{-1}=z^{-1}xz,&
 aya^{-1}=c(x,z)y c(x,z)^{-1},&
 aza^{-1}=c(x,z)z,\\
 bxb^{-1}=x,&
 byb^{-1}=z^{-1}yz,&
 bzb^{-1}=c(y,z)z.
\end{array}                                                     \tag{FCQ.3}
\]
They are the actual presentation used for the geometric reconstruction
in B.5.  The displayed braid matrices satisfy the braid relations,
and $z_4$ maps to $I$, so $q_3$ is a defined homomorphism on $P$.
The breadth-first enumeration using the five displayed generators
visits all $768$ pairs $(\ell,q_3)$; hence $q$ is onto.

The cellular chain complex of the presentation cover has one vertex
for each element of $Q$, five oriented edges at each vertex, and six
relator cells at each vertex.  Its first homology is exactly
$N/[N,N]N^3$.  This last assertion uses only the presentation and the
covering-space identification of its fundamental group with $N$;
it does not require an asphericity assertion about its two-complex.

For a small exact matrix model, Fourier-decompose the regular
$\mathbb F_3[\mathbb F_2^5]$ module.  Its $32$ characters have values
$\{1,-1\}\subset\mathbb F_3$, and their idempotents exist since
$32$ is invertible in $\mathbb F_3$.  For the character indexed by
$e\in\{0,\ldots,31\}$, with bit pairing $e\cdot v$, the chain
spaces have dimensions
\[
 C_2(e)=\mathbb F_3^{144},\qquad
 C_1(e)=\mathbb F_3^{120},\qquad
 C_0(e)=\mathbb F_3^{24}.
\]
Order $\operatorname{SL}_2(\mathbb F_3)$ by breadth-first enumeration
from $I$, with the ordered five generators above.  An edge is $(g,j)$,
with $0\leq g<24$ and $1\leq j\leq5$.  Its boundary is
\[
 \partial_e(g,j)=-[g]+(-1)^{e_j}[gq_3(g_j)].       \tag{FCQ.4}
\]
A relator row is its literal lifted edge chain.  At a prefix with
parity $v$, a positive $j$ contributes $(-1)^{e\cdot v}$ to the
edge at that prefix, and a negative $j$ contributes the negative
of that coefficient at the endpoint after traversing the inverse
edge.  Thus every entry of both boundary matrices is specified by
\textup{(FCQ.3)--(FCQ.4)}, without an imported representation.

Row reduction gives the following complete list of first-homology
dimensions, in character order $0,1,\ldots,31$:
\[
\begin{split}
 &(18,1,1,1,1,1,12,0,1,12,1,0,1,0,0,12,\\
 &\hspace{28mm}12,1,1,0,1,0,12,0,1,12,0,0,0,0,0,12).
\end{split}                                                     \tag{FCQ.5}
\]
Their sum is $114$.  In each block the program computes a matrix
$Z_e$ of cycle representatives and a coordinate projection $P_e$
such that $Z_e\partial_e=0$, $Z_eP_e=I$, and all relator rows are
annihilated by $P_e$.  These matrices, including zero-dimensional
blocks, are saved in the certificate.  This proves the stated
order of $R$.

\paragraph{The actual cubic action.}
Direct multiplication in the faithful Artin action of $B_4$ on
$F_4$ verifies $d^3=z_4$.  To compute its action on the five pure
generators, conjugate their literal braid words by $d$, subtract the
unique power of $z_4$ making the $A_{12}$ linking number zero, and
comb the remaining pure braid into $(a,b,x,y,z)$.  Each resulting
word is checked again by the faithful Artin action.  This provides
an exact automorphism $\theta=\operatorname{Ad}(d)$ of $P$, not
merely an automorphism of its homology.

On $Q$ its action is
\[
 (v,g)\longmapsto (D_2v,D_3gD_3^{-1}),\qquad
 D_3=U^2LU=\begin{pmatrix}2&1\\2&0\end{pmatrix},  \tag{FCQ.6}
\]
where the five columns of $D_2$ are the exponent vectors modulo two
of those five combed words.  The script saves all the words and
$D_2$, and checks \textup{(FCQ.6)} on every generator.  Thus $N$ is
preserved.  The induced edge map substitutes the same combed words
at the vertex $D_3gD_3^{-1}$ and sends character $e$ to
$eD_2^{-1}$.  Applying $Z_e$ and $P_{eD_2^{-1}}$ gives an explicitly
specified $114\times114$ matrix $M$ on $H_1(N,\mathbb F_3)$.
The certificate verifies
\[
 M^3=I,
 \qquad \operatorname{rank}(M-I)=74.              \tag{FCQ.7}
\]
It also verifies the edge substitutions against actual Schreier
loops.  The edge-map definition and the faithful braid checks are
the reason this is the full induced homology action.

\paragraph{The complete affine lifting criterion.}
There are exactly twelve fixed points of \textup{(FCQ.6)} in $Q$.
For each such point $q_0$, let $t_{q_0}$ be its breadth-first word
in the five positive generators.  Put
\[
 v_{q_0}=[\theta(t_{q_0})t_{q_0}^{-1}]
       \in H_1(N,\mathbb F_3).
\]
Every element of $R$ over $q_0$ is uniquely $u t_{q_0}$ with
$u\in H_1(N,\mathbb F_3)$.  It centralizes $d$ precisely when
\[
                 u(M-I)+v_{q_0}=0.                \tag{FCQ.8}
\]
Compute a $114\times40$ matrix $C$ with kernel on row vectors
exactly the row space of $M-I$.  Then \textup{(FCQ.8)} is soluble
exactly when $v_{q_0}C=0$.  This is an equivalence: once that vector
is zero, ordinary linear elimination provides the required $u$.
Consequently it tests every element of $R$ over the residue, with
no search truncated inside the order-$3^{114}$ kernel.

Here is the complete finite certificate.  Encode the parity vector
by its binary integer, so $26=(0,1,0,1,1)$, and write a matrix as
its four row-major entries.  The final column gives a one-based
coordinate of $v_{q_0}C$ and its nonzero value; a dash denotes the
zero vector.  A digit in the word is an index in $(a,b,x,y,z)$.
\[
\begin{array}{r|c|l|c}
 \text{parity}&q_3(q_0)&t_{q_0}&
     (\text{coordinate},\text{value})\\\hline
0 &(1,0,0,1)&1&-\\
0 &(2,1,2,0)&11&(3,2)\\
26&(2,0,0,2)&245&(5,2)\\
0 &(0,2,1,2)&1111&(4,2)\\
0 &(0,1,2,1)&2112&(4,2)\\
26&(1,2,1,0)&11245&(3,2)\\
26&(2,1,2,0)&21514&(3,2)\\
0 &(2,0,0,2)&112112&(5,2)\\
26&(0,1,2,1)&1111245&(4,2)\\
26&(0,2,1,2)&1121514&(4,2)\\
0 &(1,2,1,0)&11112112&(3,2)\\
26&(1,0,0,1)&111121514&(36,1)
\end{array}                                                     \tag{FCQ.9}
\]
In the first row $1$ denotes the empty word.  All other words use
only their displayed generator digits.  The entire $40$-coordinate
obstruction vector for every row is retained in the JSON output.
Only the identity residue passes.  Over it the solutions are
$\ker(M-I)$, of dimension $40$ by \textup{(FCQ.7)}.  This proves
\textup{(FCQ.1)} and the claimed order of the centralizer.

For the particularly relevant last row, the saved obstruction
functional has only five nonzero entries in the last character
block, which has dimension twelve.  In zero-based block coordinates
it is
\[
                  (2,1,2,1,0,1,0,0,0,0,0,0).
\]
It annihilates $M-I$ and evaluates to $1$ on
$[\theta( a^4 b a z a y)(a^4 b a z a y)^{-1}]$.
Thus the exclusion of \textup{(FCQ.2)} also has a single explicit
linear obstruction, independently readable from the complete table.
\end{proof}

\paragraph{Files and the exact scope of this component.}
The construction and every table entry are generated by
\texttt{code/check\_cubic\_centralizer\_homology.py};
its \texttt{.out}, \texttt{.json}, and \texttt{.npz} files contain,
respectively, the execution transcript, all residue words and
obstruction vectors, and all cycle/projection/action matrices.
Its finite-cover input is independently generated by
\texttt{build\_fivepoint\_cm\_level.py} and
\texttt{test\_cm\_cubic\_refinement.py}.  These scripts use exact
integer arithmetic and reduction modulo two or three.

Equation \textup{(FCQ.1)} makes the $Q$-image of a cubic conjugator
unique whenever its conjugacy equation is soluble in $R$: two such
conjugators differ by $C_R(d)$.  It supplies a finite cubic detector.
The equality of its parity with the harmonic frame parity requires
the additional comparison in A.4; \textup{(FCQ.1)} does not by itself
compute every genuine gentle-group image or finish that comparison.
